\chapter{Tangent bundles}
\label{chap:tangent}

The tangent bundle and the tangent map of a manifold modeled on a locally convex space, the
double tangent bundle with its canonical flip, and curves with their velocities
\cite[\S3.3]{GN26}, \cite[Ch.~III]{Lan99}. As in Chapter~\ref{chap:manifolds}, no completeness
is assumed.

\section{The tangent bundle and tangent maps}

\begin{definition}[The tangent bundle]
  \label{def:tangent-bundle}
  \lean{LocallyConvex.Manifolds.tangentCoordChange, LocallyConvex.Manifolds.tangentBundleCore,
    LocallyConvex.Manifolds.TangentBundle, LocallyConvex.Manifolds.tangentChart}
  \leanok
  \uses{def:mb-manifold}
  For charts $i$, $j$ of the atlas of $M$ and $x$ in both domains, the \emph{coordinate change}
  of tangent vectors is $v \mapsto \dd(j \circ i^{-1})(i(x))\,v$. The \emph{tangent bundle}
  $\T M$ is the fiber bundle with fiber $E$ obtained by gluing the trivial bundles over the
  chart domains along these coordinate changes. For a chart $i$, the map
  $\T M \supseteq \T M|_{\operatorname{dom} i} \to E \times E$ that sends $(x, v)$ to
  $(i(x), v_i)$, with $v_i$ the coordinates of $v$ in the chart $i$, is the \emph{chart of $\T M$
  over $i$}.
\end{definition}

In Lean, $\T M$ is the total space of a Mathlib \texttt{FiberBundleCore} indexed by the charts of
the atlas, whose index at $x$ is the preferred chart $c_x$. Thus a point of $\T M$ is a pair
$(x, v)$ with $x \in M$ and $v \in E$, and \emph{$v$ is the tangent vector written in the
coordinates of the preferred chart at $x$}. In any other chart $i$, its coordinates are
$\dd(i \circ c_x^{-1})(c_x(x))\,v$. This is a choice of representative, not a restriction:
every construction below (tangent maps, local representatives of vector fields, Lie brackets,
derivations, velocities) is defined in the preferred chart and then shown to transform correctly
in every chart of the atlas. The coordinate changes satisfy the cocycle condition by the chain
rule (\texttt{LocallyConvex.Manifolds.tangentCoordChange\_comp}), and they are continuous in $(x, v)$ since
the transition maps are $\Ck{1}$.

\begin{theorem}[$\T M$ is a manifold]
  \label{thm:man-tangent-manifold}
  \lean{LocallyConvex.Manifolds.smoothMBOn_tangentChart_transition, LocallyConvex.Manifolds.tangentChart_transition_mem,
    LocallyConvex.Manifolds.tangentChart_mem_atlas, LocallyConvex.Manifolds.tangentChart_transition_apply}
  \leanok
  \uses{def:tangent-bundle}
  The charts of $\T M$ over the charts of $M$ form an atlas modeled on $E \times E$ whose
  transition maps are
  \[
    (y, w) \longmapsto \bigl(\tau(y), \dd\tau(y)\,w\bigr), \qquad \tau = j \circ i^{-1}.
  \]
  In particular they are smooth, and $\T M$ is a manifold modeled on $E \times E$.
\end{theorem}

\begin{proof}
  \leanok
  The formula is the cocycle property of the coordinate changes. The map $(y, w) \mapsto
  \dd\tau(y)\,w$ is smooth because $\tau$ is, by the recursive definition of $\Ck{k}$, and pairs of
  smooth maps are smooth. The inverse transition has the same form with $\tau^{-1}$.
\end{proof}

\begin{definition}[The tangent map]
  \label{def:tangent-map}
  \lean{LocallyConvex.Manifolds.tangentMap}
  \leanok
  \uses{def:tangent-bundle, def:contmdiff}
  Let $f \colon M \to N$. The \emph{tangent map} $\T f \colon \T M \to \T N$ sends $(x, v)$ to
  \[
    \bigl(f(x),\ \dd(c'_{f(x)} \circ f \circ c_x^{-1})(c_x(x))\,v\bigr),
  \]
  where $c_x$ and $c'_{f(x)}$ are the preferred charts at $x$ and $f(x)$.
\end{definition}

The definition makes sense for every $f$, since $\dd$ is a total function in Lean; it is
meaningful when $f$ is $\Ck{1}$.

\begin{theorem}[The tangent map in charts; chain rule]
  \label{thm:man-tangent-map}
  \lean{LocallyConvex.Analysis.dMB_eq_comp_of_eventuallyEq, LocallyConvex.Manifolds.dMB_chart_comp,
    LocallyConvex.Manifolds.dMB_comp_chart, LocallyConvex.Manifolds.tangentMap_chart, LocallyConvex.Manifolds.tangentMap_id,
    LocallyConvex.Manifolds.tangentMap_comp}
  \leanok
  \uses{def:tangent-map, thm:man-contmdiff-chart}
  Let $f \colon M \to N$ be $\Ck{1}$. The derivatives of the local representatives of $f$
  transform under changes of charts by the derivatives of the transition maps. Consequently, in
  the charts of $\T M$ and $\T N$ over charts $i$ of $M$ and $j$ of $N$, the tangent map reads
  \[
    (y, w) \longmapsto \bigl(h(y), \dd h(y)\,w\bigr), \qquad h = j \circ f \circ i^{-1}.
  \]
  Moreover $\T\id = \id$ and, for $\Ck{1}$ maps $f$ and $g$, $\T(g \circ f) = \T g \circ \T f$.
\end{theorem}

\begin{proof}
  \leanok
  Near a point, the local representative of $f$ in one pair of charts agrees with the
  composition of its local representative in another pair with transition maps; the chain rule for maps that agree with a composition near the point gives
  the transformation laws. The formula in charts follows by composing these laws with the
  coordinate changes of $\T M$ and $\T N$, and functoriality is the chain rule in the preferred
  charts.
\end{proof}

\begin{theorem}[Smoothness of the tangent map]
  \label{thm:man-tangent-map-smooth}
  \lean{LocallyConvex.Manifolds.ContMDiffMB.tangentMap}
  \leanok
  \uses{thm:man-tangent-map, thm:man-tangent-manifold}
  If $f \colon M \to N$ is $\Ck{k+1}$, then $\T f \colon \T M \to \T N$ is $\Ck{k}$.
\end{theorem}

\begin{proof}
  \leanok
  In the charts of the tangent bundles $\T f$ is $(y, w) \mapsto (h(y), \dd h(y)\,w)$, and
  $\dd h$ is $\Ck{k}$ on its domain because $h$ is $\Ck{k+1}$.
\end{proof}

\section{The double tangent bundle}

The double tangent bundle is $\T\T M = \T(\T M)$, the tangent bundle of the manifold $\T M$
(Theorem~\ref{thm:man-tangent-manifold}). We follow \cite[Cap.~5]{Lar22} and \cite[\S
IV.3]{Lan99}. A chart $i$ of $M$ gives a chart of $\T M$ and then a chart of $\T\T M$, in which
points are written $((x, v), (u, w))$: $(x, v)$ is the base point in $\T M$, and $(u, w)$ the
tangent vector at it.

\begin{theorem}[Transition maps of $\T\T M$]
  \label{thm:man-double-transition}
  \lean{LocallyConvex.Manifolds.tangentAtlas, LocallyConvex.Manifolds.dMB_tangentChart_transition,
    LocallyConvex.Manifolds.DoubleTangentBundle, LocallyConvex.Manifolds.ddChart, LocallyConvex.Manifolds.ddChart_transition_apply}
  \leanok
  \uses{thm:man-tangent-manifold}
  Let $\tau = j \circ i^{-1}$. In the charts of $\T\T M$ over $i$ and $j$, the transition map is
  \[
    ((x, v), (u, w)) \longmapsto \bigl((\tau x, \dd\tau_x v),\ (\dd\tau_x u,\ \dd^2\tau_x(u, v)
    + \dd\tau_x w)\bigr).
  \]
\end{theorem}

\begin{proof}
  \leanok
  The transition maps of $\T\T M$ are those of the tangent bundle of $\T M$, namely the
  transition map $\sigma(x, v) = (\tau x, \dd\tau_x v)$ of $\T M$ together with its derivative.
  Differentiating $\sigma$ in the direction $(u, w)$ gives $(\dd\tau_x u, \dd^2\tau_x(u, v) +
  \dd\tau_x w)$, since $\dd\tau_x$ is linear in $v$.
\end{proof}

\begin{theorem}[The canonical flip]
  \label{thm:man-flip}
  \lean{LocallyConvex.Manifolds.swapUV, LocallyConvex.Manifolds.ddChart_transition_swap, LocallyConvex.Manifolds.ddChartAt,
    LocallyConvex.Manifolds.flip, LocallyConvex.Manifolds.flip_proj_proj, LocallyConvex.Manifolds.ddChart_flip,
    LocallyConvex.Manifolds.flip_flip}
  \leanok
  \uses{thm:man-double-transition}
  The swap $((x, v), (u, w)) \mapsto ((x, u), (v, w))$ commutes with the transition maps of
  $\T\T M$. Hence the map $\kappa \colon \T\T M \to \T\T M$ defined as the swap in the preferred
  chart reads as the swap in every chart of $\T\T M$ over a chart of $M$. It covers the identity
  of $M$, and $\kappa \circ \kappa = \id$.
\end{theorem}

\begin{proof}
  \leanok
  By Theorem~\ref{thm:man-double-transition}, the claim amounts to $\dd^2\tau_x(u, v) =
  \dd^2\tau_x(v, u)$, the symmetry of the second differential. The rest is immediate from the
  coordinate expression.
\end{proof}

The smoothness of $\kappa$ is not stated in the library.

\section{Curves and velocities}

Curves are maps $\alpha \colon \R \to M$. We treat them, like functions, through their local
representatives, without regarding $\R$ as a manifold.

\begin{definition}[$\Ck{k}$ curves and lifts]
  \label{def:ck-curve}
  \lean{LocallyConvex.Manifolds.IsCkCurve, LocallyConvex.Manifolds.liftCoord, LocallyConvex.Manifolds.IsCkLift}
  \leanok
  \uses{def:mb-manifold, def:tangent-bundle}
  A curve $\alpha \colon \R \to M$ is $\Ck{k}$ if it is continuous and $i \circ \alpha$ is
  $\Ck{k}$ on the open set $\alpha^{-1}(\operatorname{dom} i)$ for every chart $i$ of the atlas. A
  \emph{lift} of $\alpha$ is a curve $\beta \colon \R \to \T M$ with $\beta(t) \in \T_{\alpha(t)}
  M$; its fiber coordinate in the chart $i$ is the second component $\mu_i(t)$ of the chart of
  $\T M$ over $i$ at $\beta(t)$. The lift is $\Ck{k}$ if every $\mu_i$ is $\Ck{k}$ on
  $\alpha^{-1}(\operatorname{dom} i)$.
\end{definition}

\begin{definition}[Velocity]
  \label{def:velocity}
  \lean{LocallyConvex.Manifolds.velocity}
  \leanok
  \uses{def:ck-curve}
  The \emph{velocity} of a curve $\alpha$ is the lift $\dot\alpha(t) = (\alpha(t),
  (c_{\alpha(t)} \circ \alpha)'(t))$, the derivative of the local representative in the preferred
  chart at $\alpha(t)$.
\end{definition}

\begin{theorem}[The velocity in charts]
  \label{thm:man-velocity}
  \lean{LocallyConvex.Manifolds.dMB_chart_curve, LocallyConvex.Manifolds.liftCoord_velocity,
    LocallyConvex.Manifolds.isCkLift_velocity, LocallyConvex.Manifolds.continuous_velocity}
  \leanok
  \uses{def:velocity, thm:man-tangent-manifold}
  Let $\alpha$ be a $\Ck{1}$ curve. For charts $i$, $j$ around $\alpha(t)$ we have
  $(j \circ \alpha)'(t) = \dd(j \circ i^{-1})(i(\alpha(t)))\,(i \circ \alpha)'(t)$; hence the fiber
  coordinate of $\dot\alpha(t)$ in any chart $j$ is $(j \circ \alpha)'(t)$. The velocity is a
  continuous curve in $\T M$, and if $\alpha$ is $\Ck{k+1}$, then $\dot\alpha$ is a $\Ck{k}$ lift.
\end{theorem}

\begin{proof}
  \leanok
  Near $t$, $j \circ \alpha = (j \circ i^{-1}) \circ (i \circ \alpha)$, and we apply the chain
  rule. The second statement is the first one with $i$ the preferred chart. In the chart of
  $\T M$ over $i$, $\dot\alpha$ reads $s \mapsto (i(\alpha(s)), (i \circ \alpha)'(s))$, which is
  continuous (respectively $\Ck{k}$ in the second component); the charts of $\T M$ are
  homeomorphisms onto their images.
\end{proof}

\begin{proposition}[Germs and the chain rule for velocities]
  \label{prop:man-velocity-chain}
  \lean{LocallyConvex.Manifolds.velocity_congr, LocallyConvex.Manifolds.tangentMap_velocity}
  \leanok
  \uses{def:velocity, def:tangent-map}
  The velocity $\dot\alpha(t)$ only depends on the germ of $\alpha$ at $t$. If $p \colon N \to
  M$ is $\Ck{1}$ and $\Gamma$ is a $\Ck{1}$ curve in $N$, then $\T p\,(\dot\Gamma(t))$ is the velocity of $p \circ \Gamma$
  at $t$.
\end{proposition}

\begin{proof}
  \leanok
  Directional derivatives only depend on germs. For the chain rule, near $t$ we write
  $c' \circ p \circ \Gamma = (c' \circ p \circ c^{-1}) \circ (c \circ \Gamma)$ in the preferred
  charts $c$ at $\Gamma(t)$ and $c'$ at $p(\Gamma(t))$, and apply the chain rule.
\end{proof}

\begin{remark}
  \label{rem:man-curves-ode}
  The continuity of $\dot\alpha$ in $\T M$ is what makes lengths of curves meaningful for Finsler
  metrics (see Definition~\ref{def:finsler-length-structure}); the chain rule for velocities is used for quotient
  metrics, where curves in $N$ are pushed down along a projection $p$. On the other hand, we do
  not claim the existence of integral curves of vector fields, or of curves with prescribed
  velocity: this would require existence theorems for ordinary differential equations, which
  fail in general locally convex spaces, and in Banach spaces need completeness.
\end{remark}
